\documentclass[12pt,a4paper]{article}
\usepackage[utf8]{inputenc}
\usepackage[T2A]{fontenc}
\usepackage[english,russian]{babel}
\usepackage{microtype}
\usepackage{indentfirst}
\usepackage{amsmath,amssymb,amsthm}
\usepackage{geometry}
\usepackage{booktabs}
\usepackage{tabularx}
\usepackage{array}
\usepackage{hyperref}
\usepackage[raggedright]{titlesec}

\geometry{left=2.5cm,right=2.5cm,top=2.5cm,bottom=2.5cm}

\newtheorem{theorem}{Теорема}
\newtheorem{postulate}{Постулат}
\newtheorem{definition}{Определение}

\title{\textbf{Топологическая архитектура ментальной каузальности: аппарат пучков, этальных пространств и калибровочных связностей в теории бессубстратной онто-информации}}
%\author{Анатолий Панкратов}
\author{
    \textbf{Анатолий Панкратов} \\
    \small К.ф.-м.н, независимый исследователь \\
  %  \small Department of Philosophy, Your University \\
    \small Москва, Россия \\
    \small Email: \texttt{xapankratov@gmail.com}
}

\date{2026}

\begin{document}

\maketitle

\begin{abstract}
В статье предлагается формальное решение проблемы нисходящей ментальной каузальности (\textit{downward causation}), позволяющее совместить автономную волевую агентность субъекта с законом причинной замкнутости физического мира. Опираясь на онтический структурный реализм и парадигму бессубстратной онто-информации (ОИ), автор заменяет классическую силовую модель взаимодействия сознания и материи топологическим аппаратом теории пучков и расслоений. Физический хронотоп описывается как базовое пространство Каркаса Действия ($X_{\text{КД}}$), над которым развернуто этальное пространство Универсального Топоса Логического Каркаса ($E_{\text{ЛК}}$). Показано, что на микроуровне нейрофизиологических процессов вариационный функционал Принципа наименьшего действия обладает фундаментальным калибровочным вырождением ($\delta \mathcal{S} = 0$). Ноэтический макрокод квалиа ($\Psi_{\text{ноэт}}$), представляющий собой продукт предельной алгоритмической компрессии, оптимизированной под термодинамический предел Ландауэра, осуществляет безэнергетический вертикальный калибровочный сдвиг ($g \in G_{\text{ЛК}}$) в слое этального пространства. Посредством перестройки калибровочной связности Синхронного Информационного Резонанса ($\mathcal{D}_{\text{СИР}}$) и обнуления когомологических препятствий Чеха ($\check{H}^1$) волевой акт переопределяет граничные условия --- Контекстные Топологические Инварианты (КТИ), направляя каскад Синхронных Онто-Инициаций (СОИ) вдоль выбранного горизонтального лифта. Доказывается, что при данном механизме тензор энергии-импульса физического носителя остается строго дивергентно-свободным ($\nabla_\mu T^{\mu\nu} = 0$), что снимает дилемму эпифеноменализма и индетерминизма, а также преодолевает вычислительные и термодинамические ограничения моделей IIT и Orch-OR.

\medskip
\noindent\textbf{Ключевые слова:} ментальная каузальность, причинная замкнутость, теория пучков, этальное пространство, квалиа, Принцип наименьшего действия, Принцип Ландауэра, онтический структурный реализм.
\end{abstract}

\section{Введение: тупик энергетического редукционизма и аргумент каузального исключения}

Проблема ментальной каузальности остается центральным нервом современной аналитической философии сознания и метафизики науки. Попытка объяснить, каким образом субъективные феноменальные состояния (квалиа) и акты сознательной воли способны инициировать физические изменения в макромире, неизбежно наталкивается на \textbf{аргумент каузального исключения} Джегвона Кима (Kim, 2005). Этот аргумент опирается на два фундаментальных постулата естественнонаучной картины мира: принцип \textbf{причинной замкнутости физического мира} (\textit{Causal Closure of the Physical}), согласно которому каждое физическое событие имеет достаточную физическую причину, и законы сохранения энергии-импульса.

В рамках классической субстанциальной парадигмы, восходящей к картезианскому дуализму и физикализму Нового времени, любая каузация мыслится по образцу механического или полевого взаимодействия --- как силовой трансфер энергии между материальными носителями. В этой оптике исследователь оказывается зажат в тисках неразрешимой трилеммы:
\begin{enumerate}
    \item \textbf{Интеракционистский дуализм}, постулирующий прямое силовое воздействие нематериального разума на нейроны мозга, вынужден допустить локальное нарушение первого начала термодинамики и дивергенции тензора энергии-импульса ($\nabla_\mu T^{\mu\nu} \neq 0$), что противоречит всему массиву экспериментальной физики.
    \item \textbf{Редуктивный физикализм и эпифеноменализм} спасают причинную замкнутость ценой полного лишения сознания каузальной эффективности: субъективный опыт объявляется пассивной «эпистемической тенью» нейрофизиологических процессов, не способной влиять на поведение организма.
    \item \textbf{Квантово-индетерминистские модели сознания} (Eccles, 1994; Stapp, 2007; Hameroff \& Penrose, 2014) пытаются встроить ментальную агентность в зазоры квантовой неопределенности, апеллируя к объективному коллапсу волновой функции. Однако такой подход либо привязывает сознание к уязвимым биологическим субстратам (проблема теплового декогерирования), либо подменяет разумную свободу воли стохастическим произволом микромира.
\end{enumerate}

Корень указанного концептуального тупика заключается в устаревшем субстанциалистском языке, приписывающем физической материи статус первичного «носителя», а информации --- статус вторичного эпифеномена, жестко привязанного к термодинамическим затратам по принципу Ландауэра (Landauer, 1961). Опираясь на достижения онтического структурного реализма (Ladyman \& Ross, 2007; French, 2014) и радикальный эмпиризм, настоящая работа предлагает принципиально иную онтологическую стратегию: \textbf{деконструкцию силовой каузации и переход к топологической модели калибровочного управления}.

Целью данной статьи является построение строгого математико-философского формализма \textbf{несиловой ментальной каузальности} на базе теории пучков, этальных пространств и связностей (Mac Lane \& Moerdijk, 1992). Мы покажем, что физическая реальность представляет собой эмерджентное сечение пучка бессубстратной онто-информации (ОИ), развернутого над базовым пространством Каркаса Действия ($X_{\text{КД}}$) из предвечного этального пространства Универсального Топоса Логического Каркаса ($E_{\text{ЛК}}$). В рамках этой архитектуры сознание не «толкает» материю и не впрыскивает энергию в синаптические сети. Вместо этого ноэтический макрокод квалиа функционирует как \textbf{топологический калибровочный селектор}, который безэнергетически перестраивает граничные условия --- Контекстные Топологические Инварианты (КТИ) --- внутри вырожденного спектра траекторий Принципа наименьшего действия. Это позволяет субъекту направлять макроскопическую динамику тела по оптимальному руслу при строгом сохранении причинной замкнутости физического хронотопа.

\section{Онтологический базис: Бессубстратная онто-информация и регулятивные каркасы (ЛК и КД)}

\subsection{Деконструкция субстрата и каузальная асимметрия Принципа Ландауэра}
Традиционный физикализм покоится на метафизическом допущении, согласно которому любое наблюдаемое свойство или информационная структура нуждается в изолированном материальном носителе --- «субстрате» (\textit{substantia}). Однако радикальный эмпиризм современной квантовой теории поля демонстрирует избыточность этой категории: принцип тождественности (неразличимости) элементарных частиц лишает объекты микромира индивидуальной материальной «подложки», оставляя в распоряжении исследователя исключительно инварианты, симметрии и математические структуры отношений (Ladyman \& Ross, 2007; French, 2014).

Главным аргументом в пользу первичности материи над информацией обычно выступает \textbf{Принцип Ландауэра} (Landauer, 1961), фиксирующий минимальную диссипацию тепловой энергии при логически необратимом стирании или перезаписи одного бита на физическом носителе:
\begin{equation}
W_{\min} = k_{\text{B}}T \ln 2
\end{equation}
где $k_{\text{B}}$ --- постоянная Больцмана, а $T$ --- абсолютная температура системы. Из этого соотношения материалистический редукционизм выводит тезис о том, что информация энергетически детерминирована и представляет собой вторичный эпифеномен материи.

Строгий анализ выявляет здесь фундаментальную \textbf{каузальную асимметрию}: термодинамический налог Ландауэра определяет исключительно физическую стоимость переключения материального триггера (локального интерфейса), но остается полностью индифферентным к семантической архитектуре и масштабному потенциалу самой структуры. Затраты АТФ биологическим мозгом ($\approx 20\text{ Вт}$) инвариантны относительно того, удерживает ли когнитивная система компактный алгоритмический закон гравитации при моделировании столкновения галактик или репрезентирует пустой объект. Информация выступает как \textbf{безэнергетический регулятивный селектор}: она не совершает механической работы сама, но задает топологию перенаправления энергетических потоков. Это позволяет постулировать первичный онтологический статус \textbf{бессубстратной онто-информации (ОИ)}.

\subsection{Двухуровневая архитектура регулятивных каркасов}
Чтобы объяснить, каким образом безэнергетическая онто-информация порождает стабильный, динамический хронотоп, избегая «парадокса статического кристалла» математического платонизма (Tegmark, 2014), вводится двухуровневая архитектура регулятивных матриц:
\begin{enumerate}
    \item \textbf{Логический Каркас (ЛК) как алетическая мета-топология.} ЛК представляет собой первичный онтологический горизонт, задающий абсолютные границы \textit{possibilia} (пространства возможных состояний) через оператор строгого запрета логического противоречия. В терминах теории категорий ЛК является \textbf{универсальным топосом}, внутри которого примитивные реляционные инварианты --- \textbf{Онто-Монады (ОМ)} --- и их непротиворечивые конгломераты --- \textbf{Свободные Онто-Комплексы (СОК)} --- существуют вне времени и пространства в статусе чистых пред-объектов.
    \item \textbf{Каркас Действия (КД) как вариационный оператор экстремальной оптимизации.} Физическая Вселенная является строго ограниченным эмерджентным подмножеством универсального топоса ЛК. Выделение этого подмножества осуществляет \textbf{Каркас Действия (КД)}, функционирующий как тотальный онтологический вариационный оператор на основе обобщенного \textbf{Принципа наименьшего действия (ПНД)} Гамильтона --- Лагранжа. Оперируя на до-физическом уровне, КД минимизирует не только механический функционал действия $\mathcal{S}$, но и алгоритмическую (колмогоровскую) сложность информационных конфигураций (Li \& Vitányi, 2008).
\end{enumerate}

В данной онтологии \textbf{масса и энергия} не являются субстанциями: они представляют собой информационные инварианты, выражающие меру внутренней связности и устойчивости отношений внутри пучка ОИ, прошедшего сквозь вариационный фильтр КД.

\section{Теоретико-пучковая модель материализации: СОК, КТИ, СИР и СОИ}

\subsection{Пул предпучков и функтор пучкования реальности}
Для формализации перехода от внепространственной потенциальности топоса ЛК к локализованной структуре физического хронотопа $X_{\text{КД}}$ используется математический аппарат \textbf{теории пучков} (Mac Lane \& Moerdijk, 1992; Rosiak, 2022).

На до-физическом уровне пул непротиворечивых Свободных Онто-Комплексов в топосе ЛК образует \textbf{предпучок} $\mathcal{P}$ над топологическим пространством матрицы КД: на этом уровне структуры обладают избыточными степенями свободы и не локализованы в единую пространственную картину. Превращение предпучка $\mathcal{P}$ в когерентный \textbf{пучок бессубстратной онто-информации} $\mathcal{F}$ осуществляется \textbf{функтором пучкования (\textit{sheafification})} $a: \mathbf{PSh}(X_{\text{КД}}) \to \mathbf{Sh}(X_{\text{КД}})$, роль которого онтологически исполняет вариационный фильтр Каркаса Действия.

Процесс локального развертывания пучка $\mathcal{F}$ в физическом хронотопе описывается четырехэтапной цепочкой:
\begin{enumerate}
    \item \textbf{СОК (Свободный Онто-Комплекс):} Секция предпучка в топосе ЛК, представляющая собой внепространственный семантический инвариант с бесконечным числом степеней свободы.
    \item \textbf{КТИ (Контекстный Топологический Инвариант):} Открытое подмножество $U_i \subset X_{\text{КД}}$ в топологическом пространстве матрицы КД, функционирующее как локальное граничное условие («информационная вакансия» в ткани поля).
    \item \textbf{СИР (Синхронный Информационный Резонанс):} Акт безэнергетического изоморфного сопоставления между внепространственным СОК и локальным КТИ, устанавливаемый мгновенно в силу отсутствия пространственной метрики в ЛК.
    \item \textbf{СОИ (Синхронная Онто-Инициация):} Финальный акт фазового сжатия (\textit{invariant compaction}), при котором оператор КД отсекает избыточные степени свободы СОК. Структура принудительно локализуется в $U_i$, формируя \textbf{локальное сечение пучка} $s \in \mathcal{F}(U_i)$ (локальный узел КД), манифестирующее себя через измеримые параметры массы, заряда, спина и координат.
\end{enumerate}

\subsection{Аксиомы склейки и онтологический статус макрообъектов (ЭОК)}
В рамках теории пучков любой устойчивый макрообъект дефинируется как \textbf{Эмерджентный Онто-Комплекс (ЭОК)} --- глобальное когерентное сечение $\sigma \in \Gamma(X_{\text{КД}}, \mathcal{F})$, собранное из локальных узлов СОИ на открытом покрытии $\mathcal{U} = \{U_i\}$ региона КД. Существование и стабильность ЭОК гарантируются выполнением двух фундаментальных аксиом пучка:
\begin{itemize}
    \item \textbf{Аксиома локальной тождественности (Отделимость):} Если два макроскопических ЭОК $\sigma_1, \sigma_2 \in \Gamma(U, \mathcal{F})$ распадаются на тождественные локальные узлы СОИ на всех элементах покрытия $\{U_i\}$ ($s_{1}\vert_{U_i} = s_{2}\vert_{U_i}$), то $\sigma_1 \equiv \sigma_2$.
    \item \textbf{Аксиома когерентной сборки (Склейка):} Если в каждом локальном КТИ $U_i$ инициирован узел $s_i \in \mathcal{F}(U_i)$, и на всех непустых пересечениях $U_i \cap U_j \neq \varnothing$ выполняется условие калибровочного резонанса ($s_i\vert_{U_i \cap U_j} = s_j\vert_{U_i \cap U_j}$), то в матрице КД с необходимостью существует единственное глобальное сечение $\sigma \in \Gamma(U, \mathcal{F})$.
\end{itemize}

\section{Этальное пространство ЛК и калибровочное вырождение действия}

\subsection{Архитектура расслоения: Этальное пространство ЛК над базой КД}
Для строгого описания интерфейса между безвременным миром смыслов и локализованной нейрофизиологической динамикой введем конструкцию \textbf{этального пространства (\textit{espace étalé})} над топологическим пространством КД. Пусть $X_{\text{КД}}$ --- базовое топологическое пространство, открытыми множествами $U_i \subset X_{\text{КД}}$ которого выступают КТИ в синаптической и квантовой матрице мозга, а $\mathcal{F}$ --- пучок бессубстратной онто-информации на $X_{\text{КД}}$. Построим этальное пространство $E_{\text{ЛК}}$ как дизъюнктное объединение всех слоев (ростков) пучка над точками базы:
\begin{equation}
E_{\text{ЛК}} = \bigsqcup_{x \in X_{\text{КД}}} \mathcal{F}_x
\end{equation}

В данной формализации:
\begin{enumerate}
    \item \textbf{Тотальное пространство $E_{\text{ЛК}}$} эксплицирует \textbf{Универсальный Топос Логического Каркаса (ЛК)}. В каждом вертикальном слое $\mathcal{F}_x$ предвечно и безэнергетически сосуществуют все логически непротиворечивые СОК, соответствующие данному контексту.
    \item \textbf{Каноническая проекция $p: E_{\text{ЛК}} \to X_{\text{КД}}$}, являющаяся локальным гомеоморфизмом, представляет собой оператор \textbf{Синхронной Онто-Инициации (СОИ)}. Проекция $p$ осуществляет инвариантное фазовое сжатие вертикального ростка смыслов в локальные физические параметры массы, заряда и импульса.
    \item \textbf{Непрерывное локальное сечение $s: U \to E_{\text{ЛК}}$} (где $p \circ s = \text{id}_U$) соответствует локальному \textbf{узлу КД}, а глобальное когерентное сечение $\sigma \in \Gamma(X_{\text{КД}}, \mathcal{F})$ дефинирует устойчивый \textbf{Эмерджентный Онто-Комплекс (ЭОК)}.
\end{enumerate}

\subsection{Калибровочное вырождение Принципа наименьшего действия и роль Квалиа}
В открытой термодинамической системе биологического мозга, жестко ограниченной метаболическим пределом Ландауэра ($W_{\min} = k_{\text{B}}T \ln 2$), вариационный оператор КД вычисляет функционал действия $\mathcal{S}$:
\begin{equation}
\mathcal{S}[s] = \int_{U} \mathcal{L}(s, \partial s) \, d\mu_{\text{КД}}
\end{equation}

На микроуровне квантовых процессов в синаптических мембранах (например, при туннелировании ионов $\text{Ca}^{2+}$) функционал действия обладает \textbf{фундаментальным вырождением}: существует целое семейство локальных сечений $\{s_\alpha\} \subset \Gamma(U, \mathcal{F})$, для которых вариация действия тождественно равна нулю ($\delta \mathcal{S}[s_\alpha] = 0$) при строго инвариантном балансе массы-энергии.

На уровне базового пространства $X_{\text{КД}}$ эти сечения термодинамически неразличимы. В этальном пространстве $E_{\text{ЛК}}$ это вырождение выражается как наличие множества параллельных топологических листов, накрывающих один и тот же регион $U$.

Здесь вступает в действие \textbf{ноэтический макрокод квалиа} ($\Psi_{\text{ноэт}}$). Квалиа представляют собой продукт предельной алгоритмической компрессии (функтор факторизации), сворачивающий экспоненциальную мерность гильбертова пространства среды в монолитный, каузально эффективный интерфейсный вектор. Волевой акт субъекта не совершает механической работы в $X_{\text{КД}}$; он представляет собой \textbf{вертикальный калибровочный сдвиг $g \in G_{\text{ЛК}}$ внутри слоя $\mathcal{F}_x$}, переводящий фокус резонанса с одного листа этального пространства на другой:
\begin{equation}
s_{\beta}(x) = g(x) \cdot s_{\alpha}(x), \quad \text{при условии } \mathcal{S}[s_{\beta}] \equiv \mathcal{S}[s_{\alpha}]
\end{equation}

Поскольку преобразование $g$ действует строго вдоль вертикального слоя универсалий в ЛК и не изменяет интегральный инвариант энергии-импульса в базе, оно является \textbf{абсолютно безэнергетическим} и не генерирует диссипативного тепла, позволяя сознанию обходить термодинамический барьер Ландауэра.

\section{Связность СИР, когомологии Чеха и постулат сохранения причинной замкнутости}

\subsection{Перестройка связности и параллельный перенос воли}
Чтобы вертикальный выбор в этальном пространстве $E_{\text{ЛК}}$ привел к макроскопическому движению в физическом хронотопе, требуется согласование локальных узлов на всем протяжении пространственно-временного покрытия $\mathcal{U} = \{U_i\}$. В аппарате расслоений эту функцию выполняет \textbf{калибровочная связность $\mathcal{D}_{\text{СИР}}$}, онтологическим эквивалентом которой выступает \textbf{Синхронный Информационный Резонанс (СИР)}.

Механизм обратной ментальной каузальности разворачивается в три формальных шага:
\begin{enumerate}
    \item \textbf{Перестройка связности (Топологическая селекция КТИ):} Ноэтический сдвиг $g$ модифицирует форму связности $\mathcal{D}_{\text{СИР}} \to \mathcal{D}'_{\text{СИР}}$ на расслоении. По аналогии с несиловым эффектом Ааронова --- Бома (Aharonov \& Bohm, 1959), новая связность переопределяет конфигурацию граничных условий (КТИ) в матрице $X_{\text{КД}}$.
    \item \textbf{Аннулирование когомологического препятствия:} Для произвольного набора локальных сечений попытка их объединения блокируется ненулевой первой группой когомологий Чеха $\check{H}^1(\mathcal{U}, \mathcal{F}) \neq 0$, фиксирующей топологическую несовместимость соседних узлов. Связность $\mathcal{D}'_{\text{СИР}}$, индуцированная волевым актом, перестраивает коциклы на пересечениях $U_i \cap U_j$ так, что препятствие к склейке вдоль выбранного листа обнуляется.
    \item \textbf{Горизонтальный лифт (Каскад СОИ):} В обновленной топологии КТИ итерационный движок Каркаса Действия осуществляет \textbf{параллельный перенос (горизонтальный лифт)} состояния вдоль выбранного сечения $s_\beta$. Каждая последующая \textbf{Синхронная Онто-Инициация (СОИ)} квантовых событий в синапсах автоматически локализуется в новых КТИ.
\end{enumerate}

\subsection{Постулат сохранения причинной замкнутости}
Из предложенного формализма вытекает строгое онтологическое следствие:

\begin{postulate}[Инвариантность каузального трека]
Любой несиловой волевой акт, реализованный как калибровочный выбор горизонтального лифта $\sigma_\beta \in \Gamma(X_{\text{КД}}, \mathcal{F})$ в этальном пространстве $E_{\text{ЛК}}$, оставляет тензор энергии-импульса $T^{\mu\nu}$ в базовом пространстве $X_{\text{КД}}$ строго дивергентно-свободным:
\begin{equation}
\nabla_\mu T^{\mu\nu} = 0
\end{equation}
\end{postulate}

Для внешнего физического наблюдателя, регистрирующего исключительно локальные проекции $p(\sigma_\beta(x))$ в базе КД, поведение системы выглядит как непрерывная, полностью детерминированная цепь физических событий, подчиняющаяся уравнениям Эйлера --- Лагранжа. Свобода воли не нарушает детерминизм природы и не требует апелляции к индетерминизму квантового хаоса. Она укоренена в \textbf{структурной неполноте локального описания}: физический закон ПНД задает класс эквивалентных траекторий (калибровочную орбиту), тогда как выбор конкретного представителя на этой орбите принадлежит высокоуровневому когнитивному ЭОК.

\section{Компаративный анализ: пучковая модель против IIT и Orch-OR}

Предложенный аппарат несиловой калибровочной каузальности преодолевает концептуальные и термодинамические тупики двух наиболее влиятельных современных теорий сознания --- \textbf{Теории интегрированной информации (IIT)} Джулио Тонони (Tononi et al., 2016) и \textbf{модели объективной редукции (Orch-OR)} Роджера Пенроуза и Стюарта Хамероффа (Hameroff \& Penrose, 2014).

\begin{table}[htbp]
\centering
\caption{Сравнительный анализ моделей ментальной каузальности и сознания}
\label{tab:comparison}
\begin{tabularx}{\textwidth}{
  >{\bfseries\raggedright\arraybackslash\hspace{0pt}}X
  >{\raggedright\arraybackslash\hspace{0pt}}X
  >{\raggedright\arraybackslash\hspace{0pt}}X
  >{\raggedright\arraybackslash\hspace{0pt}}X
}
\toprule
Критерий сравнения & Теория IIT (Тонони) & Модель Orch-OR (Пенроуз --- Хамерофф) & Пучковая теория бессубстратной ОИ \\
\midrule
Онтологический базис опыта & Причинная структура физического субстрата (максимум $\Phi$). & Гравитационный коллапс волновой функции в микротрубочках. & Безэнергетический ноэтический макрокод $\Psi_{\text{ноэт}}$ в $E_{\text{ЛК}}$. \\
\addlinespace
Механизм ментальной каузальности & Отсутствует или редуцируется к синтаксису (риск эпифеноменализма). & Неалгоритмический объективный коллапс ($\text{OR}$). & Калибровочный выбор горизонтального лифта ($g \in G_{\text{ЛК}}$) при $\delta \mathcal{S} = 0$. \\
\addlinespace
Отношение к пределу Ландауэра & Игнорирует энтропийную стоимость вычисления $\Phi$ в реальном времени. & Требует защиты квантовой когерентности от теплового шума мозга. & Обходит предел Ландауэра за счет алгоритмической компрессии в квалиа. \\
\addlinespace
Причинная замкнутость & Сохраняется ценой утраты волевой агентности субъекта. & Потенциально нарушается невычислимым «впрыском» редукции. & Строго сохраняется ($\nabla_\mu T^{\mu\nu} = 0$) за счет сдвига КТИ. \\
\bottomrule
\end{tabularx}
\end{table}

\begin{enumerate}
    \item \textbf{Преодоление вычислительного паралича IIT:} Прямое вычисление пространства состояний даже для небольшой нейронной сети требует экспоненциальных затрат энергии, что несовместимо с метаболическим лимитом мозга ($\approx 20\text{ Вт}$) по Принципу Ландауэра. Пучковая теория ОИ доказывает, что мозг не вычисляет синтаксическую матрицу $\Phi$; квалиа $\Psi_{\text{ноэт}}$ представляют собой \textbf{ультра-сжатый функтор факторизации}, сворачивающий экспоненциальную мерность в готовый интерфейсный маркер без энергозатрат.
    \item \textbf{Избавление от «проблемы теплого мозга» и случайности в Orch-OR:} В теории бессубстратной ОИ нет необходимости «удерживать» хрупкую когерентность внутри горячего биологического носителя: нелокальная связность изначально обеспечена на уровне внепространственного этального пространства $E_{\text{ЛК}}$. Кроме того, воля реализуется через \textbf{гладкий параллельный перенос (горизонтальный лифт)} вдоль калибровочной орбиты ПНД, исключая квантовый произвол.
\end{enumerate}

\section{Заключение: свобода воли как топологический отбор в алгебраической структуре бытия}

Проведенное исследование демонстрирует, что классический конфликт между автономией сознательной воли и законом причинной замкнутости физического мира является артефактом субстанциалистской метафизики, отождествляющей любую форму каузации с силовым энергетическим обменом. Отказ от ненаблюдаемого материального субстрата в пользу онтического структурного реализма и бессубстратной онто-информации (ОИ) позволяет перевести проблему ментальной каузальности из плоскости термодинамических нарушений в область калибровочной топологии.

Построенный в работе формализм этального пространства Универсального Топоса Логического Каркаса ($E_{\text{ЛК}}$), развернутого над базовым пространством матрицы Каркаса Действия ($X_{\text{КД}}$), раскрывает строгий механизм нисходящего управления макросистемой. Доказанный постулат инвариантности каузального трека ($\nabla_\mu T^{\mu\nu} = 0$) гарантирует, что для любого внешнего физического наблюдателя динамика носителя остается строго детерминированной и энергетически замкнутой. Тем самым свобода воли получает строгое математическое обоснование как способность высокоуровневого когнитивного Эмерджентного Онто-Комплекса осуществлять безэнергетический топологический отбор в предвечном пространстве логических возможностей.

\begin{thebibliography}{99}

\bibitem{Aharonov1959}
Aharonov, Y., \& Bohm, D. (1959). Significance of Electromagnetic Potentials in the Quantum Theory. \textit{Physical Review}, 115(3), 485--491.

\bibitem{Eccles1994}
Eccles, J. C. (1994). \textit{How the Self Controls Its Brain}. Springer.

\bibitem{French2014}
French, S. (2014). \textit{The Structure of the World: Metaphysics and Representation}. Oxford University Press.

\bibitem{Hameroff2014}
Hameroff, S., \& Penrose, R. (2014). Consciousness in the universe: A review of the 'Orch OR' theory. \textit{Physics of Life Reviews}, 11(1), 39--78.

\bibitem{Kim2005}
Kim, J. (2005). \textit{Physicalism, or Something Near Enough}. Princeton University Press.

\bibitem{Ladyman2007}
Ladyman, J., \& Ross, D. (2007). \textit{Every Thing Must Go: Metaphysics Naturalized}. Oxford University Press.

\bibitem{Landauer1961}
Landauer, R. (1961). Irreversibility and Heat Generation in the Computing Process. \textit{IBM Journal of Research and Development}, 5(3), 183--191.

\bibitem{Li2008}
Li, M., \& Vitányi, P. (2008). \textit{An Introduction to Kolmogorov Complexity and Its Applications}. Springer.

\bibitem{MacLane1992}
Mac Lane, S., \& Moerdijk, I. (1992). \textit{Sheaves in Geometry and Logic: A First Introduction to Topos Theory}. Springer.

\bibitem{Rosiak2022}
Rosiak, D. (2022). \textit{Sheaf Theory through Examples}. MIT Press.

\bibitem{Stapp2007}
Stapp, H. P. (2007). \textit{Mindful Universe: Quantum Mechanics and the Participating Observer}. Springer.

\bibitem{Tegmark2014}
Tegmark, M. (2014). \textit{Our Mathematical Universe: My Quest for the Ultimate Nature of Reality}. Knopf.

\bibitem{Tononi2016}
Tononi, G., Boly, M., Massimini, M., \& Koch, C. (2016). Integrated information theory: from consciousness to its physical substrate. \textit{Nature Reviews Neuroscience}, 17(7), 450--461.

\end{thebibliography}

Данная работа распространяется под лицензией Creative Commons Attribution 4.0 International License (CC-BY 4.0).

\end{document}
